\documentclass[12pt,letterpaper]{article}
\usepackage[margin=1in]{geometry}
\usepackage[T1]{fontenc}
\usepackage{amsmath,amssymb}
\usepackage{newtxtext,newtxmath}
\usepackage{helvet}
\usepackage{enumitem}
\usepackage{caption,booktabs}
\usepackage[hidelinks]{hyperref}
\captionsetup{font=small,labelfont=bf,skip=4pt}
\renewcommand{\arraystretch}{1.15}
\setlength{\emergencystretch}{2em}
\pagestyle{empty}
\setlength{\parindent}{0pt}
\setlength{\parskip}{0pt}
\AtBeginDocument{\fontsize{12}{13.8}\selectfont}
\setlist[enumerate,1]{label=\arabic*.,leftmargin=0.25in,
labelsep=0.08in,align=left,topsep=0pt,itemsep=14pt,parsep=6pt}
\newcommand{\um}{\,\mu\mathrm{m}}
\newcommand{\uA}{\,\mu\mathrm{A}}
\newcommand{\uS}{\,\mu\mathrm{S}}
\newcommand{\mS}{\,\mathrm{mS}}
\newcommand{\kohm}{\,\mathrm{k}\Omega}
\newcommand{\Mohm}{\,\mathrm{M}\Omega}
\newcommand{\V}{\,\mathrm{V}}
\newcommand{\MHz}{\,\mathrm{MHz}}
\newcommand{\kHz}{\,\mathrm{kHz}}
\newcommand{\pF}{\,\mathrm{pF}}

\hypersetup{pdfauthor={Anonymous contributor},pdftitle={ECEN 474/704 - Lab 7 Prelab},pdfsubject={Anonymized laboratory prelab}}
\begin{document}
{\fontsize{11}{12.65}\selectfont
\hfill October 16, 2026\par}
{\centering
{\fontsize{14}{16.1}\selectfont\sffamily\bfseries ECEN 474/704\par}
{\fontsize{14}{16.1}\selectfont Lab 7: Prelab\par}
{\fontsize{14}{16.1}\selectfont Anonymous contributor\par}
}
\vspace{12pt}
\begin{enumerate}

\item \textbf{Comparison of inverting-amplifier topologies.}
Rank the four amplifiers in gain, input impedance, output impedance,
and linearity. Give expressions for the first three quantities.

\textbf{Solution.}

The comparison assumes comparable bias currents and device parameters.
The ratings favor high voltage gain, high input impedance, and low
output impedance for driving a load.

\begin{center}
\small
\setlength{\tabcolsep}{5pt}
\begin{tabular}{@{}lcccc@{}}
\toprule
Specification & Current mirror & Digital CMOS & PMOS buffer & CMOS buffer\\
\midrule
Voltage gain & Medium & Good & Poor & Poor\\
Input impedance & Good & Good & Good & Good\\
Output impedance & Poor & Poor & Good & Good\\
Linearity & Medium & Poor & Good & Medium\\
\bottomrule
\end{tabular}
\end{center}

The buffers have a target gain of $-1$ and are linear in the ideal
square-law model. Tracking between devices of the same type favors
the PMOS pair, although body effect and other nonidealities limit
its linearity.

\textbf{(a) Gain and output resistance.}
Let $g_{oi}=1/r_{oi}$ and $g_o=g_{o1}+g_{o2}$.
With the supplies and bias-reference node at AC ground,
\begin{center}
\small
\setlength{\tabcolsep}{7pt}
\renewcommand{\arraystretch}{1.25}
\begin{tabular}{@{}lcc@{}}
\toprule
Circuit & $A_{v0}$ & $R_{\mathrm{out}}$\\
\midrule
Current-mirror load & $-g_{m1}/g_o$ & $1/g_o$\\
Digital CMOS & $-(g_{m1}+g_{m2})/g_o$ & $1/g_o$\\
Self-biased PMOS & $-g_{m1}/(g_{m2}+g_{mb2}+g_o)$
 & $1/(g_{m2}+g_{mb2}+g_o)$\\
Self-biased CMOS & $-g_{m1}/(g_{m2}+g_o)$ & $1/(g_{m2}+g_o)$\\
\bottomrule
\end{tabular}
\end{center}
The output resistance is evaluated with the input held constant. A finite
$R_L$ adds $1/R_L$ to each denominator. The PMOS load's source
is the output while its body is fixed, giving the $g_{mb2}$ term.
The CMOS load's source and body are both fixed at $V_{DD}$.

\textbf{(b) Input impedance.}
Ignoring gate leakage, all four circuits have
$R_{\mathrm{in,DC}}\rightarrow\infty$.
At finite frequency, let $k(s)=v_o/v_A$, where $v_A$ is the gate node
after $R_S$. Then
\[
Z_{\mathrm{in},A}(s)=
\frac{1}{s\{C_A+C_C[1-k(s)]\}}
\approx\frac{1}{s[C_A+C_C(1+|A_{v0}|)]}.
\]
The source sees $R_S+Z_{\mathrm{in},A}$.
For the single-input-device circuits, $C_A\approx C_{gs1}$; digital CMOS has
$C_A\approx C_{gs1}+C_{gs2}$.
Similarly, $C_C=C_m+C_{gd1}$, with $C_{gd2}$ added for digital CMOS.
Thus the equal DC input-resistance ratings do not imply equal AC
input loading.

\newpage
\item \textbf{Miller-theorem derivation and comparison.}
Derive the generic amplifier transfer function using Miller's theorem,
compare it with the nodal result, and discuss the approximation's utility.

\textbf{Solution.}

\textbf{(a) Miller equivalent.}
Let $v_A$ be the input node after $R_A$, let $v_o$ be the output,
and define $A_0=G_mR_B>0$, so the DC voltage gain is $-A_0$.
For the capacitor $C_C$ between these nodes,
\begin{align*}
i_A&=sC_C(v_A-v_o)=sC_C[1-k(s)]v_A,\\
i_o&=sC_C(v_o-v_A)=sC_C[1-1/k(s)]v_o,
\qquad k(s)=\frac{v_o}{v_A}.
\end{align*}
These identities are exact if the frequency-dependent $k(s)$ is retained.
Replacing $k(s)$ by its low-frequency value $-A_0$ gives
\[
\boxed{C_{\mathrm{in}}=C_A+C_C(1+A_0),\qquad
C_{\mathrm{out}}=C_B+C_C(1+1/A_0).}
\]
The input divider and output transconductance stage then give
\[
\frac{v_A}{v_{\mathrm{in}}}\approx
\frac{1}{1+sR_AC_{\mathrm{in}}},\qquad
\frac{v_o}{v_A}\approx
\frac{-G_mR_B}{1+sR_BC_{\mathrm{out}}}.
\]
Thus the constant-gain Miller approximation is
\[
\boxed{H_M(s)\approx
\frac{-A_0}
{(1+sR_AC_{\mathrm{in}})(1+sR_BC_{\mathrm{out}})}.}
\]
The resulting time constants are $R_AC_{\mathrm{in}}$ and
$R_BC_{\mathrm{out}}$. Their accuracy as pole estimates depends on
the frequency variation of the inter-node gain.

\textbf{(b) Exact nodal result.}
With the controlled current $G_mv_A$ directed from the output to
AC ground, the two node equations are
\begin{align*}
\left[\frac{1}{R_A}+s(C_A+C_C)\right]v_A-sC_Cv_o
 &=\frac{v_{\mathrm{in}}}{R_A},\\
(G_m-sC_C)v_A+
\left[\frac{1}{R_B}+s(C_B+C_C)\right]v_o&=0.
\end{align*}
Eliminating $v_A$ yields
\[
\boxed{H(s)=\frac{-A_0(1-sC_C/G_m)}{1+a s+b s^2},}
\]
where
\begin{align*}
a&=R_A(C_A+C_C)+R_B(C_B+C_C)+G_mR_AR_BC_C,\\
b&=R_AR_B(C_AC_B+C_AC_C+C_BC_C).
\end{align*}
The denominator and zero agree with the handout. The overall
minus sign makes $H(0)=-G_mR_B$, consistent with an inverting amplifier.
The exact inter-node gain is
\[
k(s)=-\frac{G_m-sC_C}{1/R_B+s(C_B+C_C)}.
\]
Using this $k(s)$ in the Miller identities recovers the same nodal result.

\newpage
\textbf{(c) Comparison and utility.}
Expanding $H_M$ gives coefficients
$a_M=R_AC_{\mathrm{in}}+R_BC_{\mathrm{out}}$ and
$b_M=R_AR_BC_{\mathrm{in}}C_{\mathrm{out}}$.
Their differences from the exact denominator are
\begin{align*}
a_M-a&=\frac{R_BC_C}{A_0}=\frac{C_C}{G_m},\\
b_M-b&=R_AR_B\left[
\frac{C_AC_C}{A_0}+A_0C_BC_C
+C_C^2\left(2+A_0+\frac{1}{A_0}\right)\right].
\end{align*}
The constant-gain approximation also removes the right-half-plane zero
\[
\boxed{s_z=+\frac{G_m}{C_C},\qquad
f_z=\frac{G_m}{2\pi C_C}.}
\]
Miller's theorem itself is an exact circuit identity; the loss of
accuracy comes from substituting the constant $-A_0$ for $k(s)$.
The approximation therefore provides an estimate of input loading
and an initial compensation capacitance. The second pole and phase
near unity gain require the frequency-dependent gain or the nodal result.

For $A_0\gg1$, negligible $C_A$, and a dominant Miller time constant,
\[
f_{p1}\approx\frac{1}{2\pi A_0R_AC_C},\qquad
\mathrm{GBW}\approx A_0f_{p1}
\approx\frac{1}{2\pi R_AC_C}.
\]
For $R_A=100\kohm$ and $\mathrm{GBW}=2\MHz$, this gives
\[
C_C\approx0.7958\,\mathrm{pF}.
\]
This is the total feedback capacitance, including $C_{gd}$; the external
$C_m$ is smaller by the relevant parasitic capacitance.
The full two-node expression provides the corresponding refinement.

\textbf{(d) Frequency and phase conventions.}
When $a^2\geq4b$, the two real poles are at
$s=-\omega_{p1}$ and $s=-\omega_{p2}$, where
\[
\omega_{p1}=\frac{2}{a+\sqrt{a^2-4b}},\qquad
\omega_{p2}=\frac{a+\sqrt{a^2-4b}}{2b},\qquad
f_{pi}=\frac{\omega_{pi}}{2\pi}.
\]
Define the measured bandwidth $f_{\mathrm{BW}}$ by
$|H(j2\pi f_{\mathrm{BW}})|=A_0/\sqrt{2}$.
It is approximately $f_{p1}$ only when the other pole and the zero
are sufficiently far away. The measured gain--bandwidth product is
$\mathrm{GBW}=A_0f_{\mathrm{BW}}$, while the unity-gain frequency
satisfies $|H(j2\pi f_{\mathrm{UGF}})|=1$.

After removing the DC inversion, the dynamic phase lag is
\[
\phi_{\mathrm{lag}}(f)=
\tan^{-1}\!\left(\frac{f}{f_{p1}}\right)+
\tan^{-1}\!\left(\frac{f}{f_{p2}}\right)+
\tan^{-1}\!\left(\frac{f}{f_z}\right).
\]
The RHP zero adds phase lag. Therefore,
$f_{p2}\geq\sqrt{3}\,\mathrm{GBW}$ alone guarantees neither the exact
UGF nor a lag below $120^\circ$ when the zero is appreciable.
For a monotone low-pass buffer with $A_0=1$, a positive-frequency
unity-gain crossing need not exist. The buffer phase is therefore
evaluated at $f_{\mathrm{BW}}$, as specified in the handout.

\newpage
\item \textbf{Design of the four inverting amplifiers.}
Design the current-mirror-loaded inverter, digital CMOS inverter,
self-biased PMOS inverter, and self-biased CMOS inverter in
Figs.~7-5, 7-6, 7-8, and 7-9, respectively.
All four use $V_{DD}=+0.9\V$, $V_{SS}=-0.9\V$,
$R_S=100\kohm$, and $\mathrm{GBW}=2\MHz$.

\textbf{Solution.}

\textbf{Common specifications and device assumptions.}
The high-gain designs require $A_0=10^{30/20}=31.6228$,
$C_L=1\pF$, and dynamic phase lag below $120^\circ$ at UGF.
The two diode-load designs require $A_0=1$ and $C_L=25\pF$.
Here $A_0=|A_{v0}|$ is positive; the voltage gain is negative.
The nominal output is selected as $V_{o,Q}=0$.

The device parameters are taken from Table~I of the Lab~1 MOS
characterization report, measured at single-finger $W/L=5/0.27\um$:
\begin{center}
\small
\begin{tabular}{@{}lcc@{}}
\toprule
Parameter & NFET & PFET\\
\midrule
$K_P=\mu C_{ox}$ ($\mu$A/V$^2$) & 215.6 & 47.57\\
$|V_{TH}|$ (V) & 0.4330 & 0.4068\\
$\lambda$ at $L_0=0.27\um$ (V$^{-1}$) & 0.1624 & 0.1945\\
$\theta$ (V$^{-1}$) & 0.9999 & 0.6558\\
\bottomrule
\end{tabular}
\end{center}
The DC sizing uses the square-law model. Channel-length modulation
is retained only through the small-signal estimate
$g_o\approx\lambda I_D$:
\[
I_D=\frac12 K_P\frac{W_{\mathrm{eff}}}{L}V_{OV}^2,
\qquad g_m=\frac{2I_D}{V_{OV}},
\qquad r_o\approx\frac{1}{\lambda I_D}.
\]
The numerical designs initially omit body effect, mobility
degradation, and DC channel-length-modulation corrections.
The listed $\theta$ values are not applied to the square-law results.
For lengths different from $L_0$, the assumed scaling is
\[
\lambda_n(L)=0.1624\frac{L_0}{L},\qquad
\lambda_p(L)=0.1945\frac{L_0}{L},
\]
while holding $K_P$ and $V_{TH}$ fixed.
This length dependence is a design approximation that has not been
verified by device characterization.

The frequency calculation assumes $C_A=0$, $C_B=C_L$, and
$C_C=C_m$. Gate, junction, and interconnect capacitances in the PDK
modify this model. The results below are analytical estimates;
their accuracy requires circuit simulation.

\textbf{Width and multiplicity convention.}
Each table gives both total effective width and width per finger:
\[
W_{\mathrm{eff}}=w\,n_f\,m.
\]
Parts~(a)--(d) derive continuous $w_{\mathrm{ana}}$ and $L_{\mathrm{ana}}$.
The later grid-constrained table gives the initial dimensions, with $w,L$
on the $0.01\um$ grid and initial $n_f=m=2$; both counts remain
adjustable. DC and AC verification apply to the selected geometry;
any subsequent dimensional change requires recalculation and simulation.

\newpage
\textbf{(a) Current-mirror-loaded inverter (Fig.~7-5).}

\textbf{Transconductance and bias current.}
M1 is the NMOS driver; M2 and M3 form a 1:1 PMOS mirror.
The handout's second-pole estimate would give
\[
g_{m1}\gtrsim 2\pi(\sqrt{3})(2\MHz)(1\pF)=21.77\uS.
\]
This bound neglects the RHP zero. A larger value,
$g_{m1}=300\uS$, is selected, with phase evaluated using
the complete two-node transfer function.
At $L=0.270\um$, $\lambda_\Sigma=\lambda_n+\lambda_p=0.3569\V^{-1}$.
The required current and overdrive are
\begin{align*}
I_D=I_{\mathrm{bias}}
&=\frac{g_{m1}}{A_0\lambda_\Sigma}
=\frac{300\times10^{-6}}{31.6228(0.3569)}
=26.581\uA,\\
R_{\mathrm{out}}&=\frac{1}{\lambda_\Sigma I_D}=105.409\kohm,\\
V_{OV1}&=\frac{2I_D}{g_{m1}}=0.177208\V.
\end{align*}
The choice $g_{m2}=g_{m3}=g_{m1}$ gives equal overdrives for all three devices.
The required aspect ratios are
\[
\left(\frac{W_{\mathrm{eff}}}{L}\right)_n
=\frac{2(26.581)}{215.6(0.177208)^2}=7.8522,
\quad
\left(\frac{W_{\mathrm{eff}}}{L}\right)_p=35.5881.
\]
\begin{center}
\small
\begin{tabular}{@{}lccccc@{}}
\toprule
Device & $w_{\mathrm{ana}}$ ($\mu$m) & $L_{\mathrm{ana}}$ ($\mu$m) & $n_f$ & $m$ & $W_{\mathrm{eff}}$ ($\mu$m)\\
\midrule
M1 (NFET) & 0.53002 & 0.27000 & 2 & 2 & 2.12008\\
M2, M3 (PFET), each & 2.40220 & 0.27000 & 2 & 2 & 9.60878\\
\bottomrule
\end{tabular}
\end{center}
The estimated gate biases are
\begin{align*}
V_{in,Q}&=V_{SS}+V_{TN}+V_{OV1}=-0.289792\V,\\
V_{G2}=V_{G3}&=V_{DD}-|V_{TP}|-V_{OV2}=0.315992\V.
\end{align*}
With $V_{o,Q}=0$, M1 and M2 each have $|V_{DS}|=0.9\V$;
M3 is diode-connected with $V_{SD3}=0.584008\V$.
All exceed $V_{OV}=0.177208\V$ under the assumed model.
The mirror current will require adjustment when unequal drain
voltages and DC channel-length modulation are included.

\textbf{Compensation and two-node response.}
The Miller estimate gives $C_C\approx0.795775\pF$.
With $f_{\mathrm{BW}}=2\MHz/31.6228=63.2456\kHz$,
the complete transfer function satisfies the half-power condition at
$\omega_B=2\pi f_{\mathrm{BW}}$:
\begin{equation}
(1-b\omega_B^2)^2+(a\omega_B)^2
=2\left[1+\left(\frac{\omega_B C_C}{G_m}\right)^2\right].
\label{eq:halfpower}
\end{equation}
This yields $C_m=C_C=0.716828\pF$ in the no-parasitic model.
The resulting estimates are
\[
A_{v0}=-31.6228,\quad f_{p1}=63.2456\kHz,\quad
f_{p2}=53.005\MHz,\quad f_z=66.608\MHz.
\]
The actual unity-gain crossing of this model is
$f_u=1.99848\MHz$, where the dynamic phase lag is
$92.07^\circ<120^\circ$. Its GBW is $2.000\MHz$.
The external $C_m$ is reduced by the intrinsic contribution of $C_{gd1}$.

\newpage
\textbf{(b) Digital CMOS inverter (Fig.~7-6).}

\textbf{Length and current selection.}
Both gates are driven by the input, giving
$G_m=g_{mn}+g_{mp}$.
The selected bias is $V_{in,Q}=V_{o,Q}=0$, with equal lengths $L_n=L_p=L$.
The overdrives are
\[
V_{OVn}=0.9-0.4330=0.4670\V,\qquad
V_{OVp}=0.9-0.4068=0.4932\V.
\]
Equal drain currents require
\[
\frac{\beta_p}{\beta_n}
=\left(\frac{V_{OVn}}{V_{OVp}}\right)^2=0.89658,
\qquad \beta_i=K_{Pi}\frac{W_{\mathrm{eff},i}}{L_i}.
\]
The unequal thresholds mean that $\beta_n=\beta_p$ does not
place the exact transition center at zero.
Using the length-scaling assumption,
\begin{align*}
A_0&=\frac{2/V_{OVn}+2/V_{OVp}}{\lambda_n(L)+\lambda_p(L)},\\
L&=L_0\,
\frac{A_0[\lambda_n(L_0)+\lambda_p(L_0)]}
{2/V_{OVn}+2/V_{OVp}}
=0.365476\um.
\end{align*}
The choice $I_D=120\uA$ places the non-dominant pole and zero
well above the required bandwidth in the assumed model.
Then
\begin{align*}
g_{mn}&=\frac{2(120\uA)}{0.4670\V}=513.919\uS,\\
g_{mp}&=\frac{2(120\uA)}{0.4932\V}=486.618\uS,\\
G_m&=1000.537\uS,\qquad
R_{\mathrm{out}}=\frac{A_0}{G_m}=31.6058\kohm.
\end{align*}

\textbf{Device sizes and operating point.}
The effective width of each transistor follows from
$W_{\mathrm{eff}}=2I_DL/(K_PV_{OV}^2)$:
\begin{center}
\small
\begin{tabular}{@{}lccccc@{}}
\toprule
Device & $w_{\mathrm{ana}}$ ($\mu$m) & $L_{\mathrm{ana}}$ ($\mu$m) & $n_f$ & $m$ & $W_{\mathrm{eff}}$ ($\mu$m)\\
\midrule
M1 (NFET) & 0.46637 & 0.36548 & 2 & 2 & 1.86547\\
M2 (PFET) & 1.89509 & 0.36548 & 2 & 2 & 7.58037\\
\bottomrule
\end{tabular}
\end{center}
At the selected bias, $V_{DS1}=V_{SD2}=0.9\V$, which exceeds
both overdrives. The NMOS body is at $V_{SS}$ and the PMOS
body at $V_{DD}$; both active devices have zero source-body bias.

\textbf{Compensation and performance.}
Solving Eq.~\eqref{eq:halfpower} with the new $G_m$ and
$R_{\mathrm{out}}$ gives $C_m=C_C=0.754672\pF$.
The resulting estimates are
\begin{align*}
A_{v0}&=-31.6228\quad(30.000\,\mathrm{dB}),\\
f_{\mathrm{BW}}&=63.2456\kHz,\qquad \mathrm{GBW}=2.000\MHz,\\
f_{p2}&=167.913\MHz,\qquad f_z=211.006\MHz,\\
f_u&=1.99895\MHz,\qquad \phi_{\mathrm{lag}}(f_u)=89.41^\circ.
\end{align*}
The analytical response satisfies the phase requirement.
In the physical circuit, $C_C=C_m+C_{gd1}+C_{gd2}$; gate and
output capacitances also affect the AC response.

\newpage
\textbf{(c) PMOS-only inverter with self-biased load (Fig.~7-8).}

\textbf{Unity-gain sizing with finite output conductances.}
M1 is the upper driver. M2 has drain and gate at $V_{SS}$,
source at $V_o$, and body at $V_{DD}$.
For the initial no-body-effect calculation at $V_{o,Q}=0$,
\[
V_{OV2}=0.9-|V_{TP}|=0.4932\V.
\]
The handout's ideal equal-size result $V_o=-V_{in}$ neglects
finite $g_o$. Retaining $g_o$ gives
\[
A_{v0}=-\frac{g_{m1}}{g_{m2}+g_{o1}+g_{o2}}.
\]
Unity gain therefore requires $G_m=g_{m1}=G_{\mathrm{out}}$, giving
$R_{\mathrm{out}}=1/G_m$ and $A_{v0}=-1$.
The simple load-dominant estimate is
\[
G_m\approx2\pi C_L\mathrm{GBW}=314.159\uS.
\]
For the initial choice $C_m=0.010\pF$, Eq.~\eqref{eq:halfpower}
determines $G_m$ at $f_{\mathrm{BW}}=2\MHz$.
For this unity-gain case,
\[
a=2R_SC_m+\frac{C_L+C_m}{G_m},\qquad
b=\frac{R_SC_LC_m}{G_m},
\]
which gives $G_m=318.387\uS$.
The realizability of this small $C_m$ depends on the PDK's
available capacitance range and the intrinsic parasitic capacitance.

The choice $L_1=L_2=0.270\um$ gives
$\lambda_\Sigma=2\lambda_p=0.3890\V^{-1}$.
Since $g_{m2}=2I_D/V_{OV2}$,
\begin{align*}
I_D&=\frac{G_m}{2/V_{OV2}+\lambda_\Sigma}=71.642\uA,\\
V_{OV1}&=\frac{2I_D}{G_m}=0.450030\V,\\
V_{in,Q}&=V_{DD}-|V_{TP}|-V_{OV1}=0.043170\V.
\end{align*}
Applying the square-law sizing equation separately to M1 and M2:
\begin{center}
\small
\begin{tabular}{@{}lccccc@{}}
\toprule
Device & $w_{\mathrm{ana}}$ ($\mu$m) & $L_{\mathrm{ana}}$ ($\mu$m) & $n_f$ & $m$ & $W_{\mathrm{eff}}$ ($\mu$m)\\
\midrule
M1 (PFET driver) & 1.00389 & 0.27000 & 2 & 2 & 4.01555\\
M2 (PFET load) & 0.83584 & 0.27000 & 2 & 2 & 3.34335\\
\bottomrule
\end{tabular}
\end{center}
The driver is slightly wider to compensate the two output
conductances. Here $g_{m2}=290.519\uS$ and
$g_{o1}+g_{o2}=27.869\uS$.
Both PMOS devices have $V_{SD}=0.9\V$, exceeding their
respective overdrives under the no-body-effect assumption.

The complete two-node calculation gives
\[
A_{v0}=-1,\quad R_{\mathrm{out}}=3.14083\kohm,\quad
f_{\mathrm{BW}}=\mathrm{GBW}=2.000\MHz.
\]
Its pole frequencies are $2.00031\MHz$ and $161.272\MHz$;
the RHP zero is at $5067.29\MHz$.
The phase lag at $f_{\mathrm{BW}}$ is $45.73^\circ$.
This unity-DC-gain low-pass model has no positive-frequency
unity-gain crossing, so its phase is evaluated at the bandwidth.

\newpage
\textbf{Body-effect correction for the PMOS load.}
The preceding PMOS-only dimensions neglect body effect.
With M2's body fixed at $V_{DD}$ and its source at zero,
$V_{BS2}=0.9\V$. The load therefore exhibits body effect,
whereas the upper driver has zero source-body bias.
Let $\Phi=2|\phi_F|$ and use the first-order relations
\begin{align*}
|V_{TP2}|&\approx |V_{TP0}|
+\gamma_p\left(\sqrt{\Phi+V_{BS2}}-\sqrt{\Phi}\right),\\
\eta&=\frac{g_{mb2}}{g_{m2}}
\approx\frac{\gamma_p}{2\sqrt{\Phi+V_{BS2}}}.
\end{align*}
The corrected quantities become
\begin{align*}
V_{OV2}&=0.9-|V_{TP2}|,\\
G_{\mathrm{out}}&=g_{m2}(1+\eta)+\lambda_\Sigma I_D,\\
I_D&=\frac{G_m}{2(1+\eta)/V_{OV2}+\lambda_\Sigma},\\
V_{OV1}&=\frac{2}{2(1+\eta)/V_{OV2}+\lambda_\Sigma},\\
\frac{\beta_1}{\beta_2}
&=\left(1+\eta+\frac{\lambda_\Sigma V_{OV2}}{2}\right)^2.
\end{align*}
Thus M1 and M2 generally require unequal sizes even with matched
PMOS process parameters. A numerical correction requires
$\gamma_p$ and $\phi_F$, which are absent from the parameter table.
The PDK operating-point values of $g_m$, $g_{mb}$, $g_{ds}$,
and threshold voltage provide the corresponding refinement.
M2's body connection remains at $V_{DD}$ in this design.

\vspace{10pt}
\textbf{(d) CMOS inverter with self-biased load (Fig.~7-9).}

M1 is the NMOS driver. M2 is the diode-connected PMOS load,
with source and body at $V_{DD}$ and drain/gate at $V_o$.
Both active source-body voltages are zero.
At $V_{o,Q}=0$, the load overdrive is $V_{OV2}=0.4932\V$.
The selected lengths are $L_1=L_2=0.540\um$, for which the assumed scaling gives
\[
\lambda_\Sigma=(0.1624+0.1945)\frac{0.270}{0.540}
=0.17845\V^{-1}.
\]
The same $C_m=0.010\pF$ and
$G_m=318.387\uS$ as in part~(c) apply because the selected
$A_0$, $R_S$, and $C_L$ are the same.
Unity gain requires $g_{m1}=g_{m2}+g_{o1}+g_{o2}$:
\begin{align*}
I_D&=\frac{318.387\times10^{-6}}{2/0.4932+0.17845}
=75.205\uA,\\
V_{OV1}&=\frac{2I_D}{G_m}=0.472411\V,\\
V_{in,Q}&=V_{SS}+V_{TN}+V_{OV1}=0.005411\V.
\end{align*}

\newpage
\textbf{CMOS self-biased inverter: dimensions and response.}

The continuous analytical dimensions are
\begin{center}
\small
\begin{tabular}{@{}lccccc@{}}
\toprule
Device & $w_{\mathrm{ana}}$ ($\mu$m) & $L_{\mathrm{ana}}$ ($\mu$m) & $n_f$ & $m$ & $W_{\mathrm{eff}}$ ($\mu$m)\\
\midrule
M1 (NFET driver) & 0.42201 & 0.54000 & 2 & 2 & 1.68803\\
M2 (PFET load) & 1.75481 & 0.54000 & 2 & 2 & 7.01925\\
\bottomrule
\end{tabular}
\end{center}
Here $g_{m2}=304.967\uS$ and
$g_{o1}+g_{o2}=13.420\uS$.
At $V_{o,Q}=0$, $V_{DS1}=V_{SD2}=0.9\V$, so both
devices are in saturation in the analytical model.
The small-signal estimates are
\begin{align*}
A_{v0}&=-1\quad(0\,\mathrm{dB}),\qquad
R_{\mathrm{out}}=3.14083\kohm,\\
f_{\mathrm{BW}}&=\mathrm{GBW}=2.000\MHz,\\
f_{p1}&=2.00031\MHz,\qquad f_{p2}=161.272\MHz,\\
f_z&=5067.29\MHz,\qquad
\phi_{\mathrm{lag}}(f_{\mathrm{BW}})=45.73^\circ.
\end{align*}
The calculated frequency response is the same as part~(c)
because the four generic-model quantities $G_m$, $R_B$, $C_B$,
and $C_C$ have the same selected values. Device
capacitances will generally make the two responses different.

\textbf{Performance with continuous analytical dimensions.}
\begin{center}
\small
\setlength{\tabcolsep}{4pt}
\begin{tabular}{@{}lrrrr@{}}
\toprule
Quantity & Mirror load & Digital CMOS & PMOS load & CMOS load\\
\midrule
$I_D$ ($\mu$A) & 26.581 & 120.000 & 71.642 & 75.205\\
$V_{in,Q}$ (V) & $-0.28979$ & 0 & 0.04317 & 0.00541\\
$C_m$ (pF) & 0.71683 & 0.75467 & 0.01000 & 0.01000\\
$|A_{v0}|$ (dB) & 30.000 & 30.000 & 0.000 & 0.000\\
$f_{\mathrm{BW}}$ (MHz) & 0.06325 & 0.06325 & 2.000 & 2.000\\
GBW (MHz) & 2.000 & 2.000 & 2.000 & 2.000\\
UGF (MHz) & 1.99848 & 1.99895 & --- & ---\\
Lag at UGF ($^\circ$) & 92.07 & 89.41 & --- & ---\\
Lag at $f_{\mathrm{BW}}$ ($^\circ$) & 45.12 & 45.04 & 45.73 & 45.73\\
\bottomrule
\end{tabular}
\end{center}
The currents above are branch currents. The mirror-loaded
inverter also draws $I_{\mathrm{bias}}$ in its reference branch.
The summary uses unrounded dimensions and the stated approximate
model, including omission of PMOS load body effect in part~(c).

\newpage
\textbf{Grid-constrained initial dimensions.}

The following initial dimensions satisfy the design grid and replace
the continuous analytical values. Both $w$ and $L$ are selected
on the $0.01\um$ grid before simulation.
All listed instances initially have $n_f=m=2$, with
$W_{\mathrm{eff}}=4w$. These initial counts remain adjustable.
\begin{center}
\small
\setlength{\tabcolsep}{5pt}
\begin{tabular}{@{}llrrrrr@{}}
\toprule
Circuit & MOS & $w$ ($\mu$m) & $L$ ($\mu$m) & $n_f$ & $m$
& $W_{\mathrm{eff}}$ ($\mu$m)\\
\midrule
Mirror load & M1 & 0.53 & 0.27 & 2 & 2 & 2.12\\
 & M2, M3, each & 2.40 & 0.27 & 2 & 2 & 9.60\\
Digital CMOS & M1 & 0.47 & 0.37 & 2 & 2 & 1.88\\
 & M2 & 1.90 & 0.37 & 2 & 2 & 7.60\\
PMOS buffer & M1 & 1.00 & 0.27 & 2 & 2 & 4.00\\
 & M2 & 0.84 & 0.27 & 2 & 2 & 3.36\\
CMOS buffer & M1 & 0.42 & 0.54 & 2 & 2 & 1.68\\
 & M2 & 1.75 & 0.54 & 2 & 2 & 7.00\\
\bottomrule
\end{tabular}
\end{center}

The operating points and small-signal parameters below correspond
to these exact dimensions. The mirror reference current is
$I_{\mathrm{bias}}=26.5812\uA$. For the other circuits, the input
bias is obtained by imposing $V_{o,Q}=0$ on the grid-constrained
dimensions. Consequently, the digital CMOS
pair no longer has exactly zero input bias.

\begin{center}
\small
\setlength{\tabcolsep}{4pt}
\begin{tabular}{@{}lrrrr@{}}
\toprule
Quantity & Mirror load & Digital CMOS & PMOS buffer & CMOS buffer\\
\midrule
$I_D$ ($\mu$A) & 26.5812 & 119.1389 & 71.9987 & 74.9987\\
$V_{in,Q}$ (V) & $-0.289789$ & $-0.000621$ & 0.041175 & 0.005889\\
$C_m$ (pF) & 0.716827 & 0.754832 & 0.010000 & 0.010000\\
$|A_{v0}|$ (dB) & 29.9998 & 30.1075 & $-0.0384$ & $-0.0088$\\
$f_{\mathrm{BW}}$ (MHz) & 0.063247 & 0.062468 & 2.009943 & 1.994614\\
GBW (MHz) & 2.000000 & 2.000000 & 2.001070 & 1.992597\\
UGF (MHz) & 1.998480 & 1.998972 & --- & ---\\
Lag at UGF ($^\circ$) & 92.065 & 89.444 & --- & ---\\
Lag at $f_{\mathrm{BW}}$ ($^\circ$) & 45.123 & 45.039 & 45.732 & 45.727\\
\bottomrule
\end{tabular}
\end{center}

These estimates retain the same model limitations: omitted body effect
and parasitic capacitances, and an unverified $\lambda(L)$ approximation.
For the grid-constrained geometry, the high-gain compensation values
follow from Eq.~\eqref{eq:halfpower} and
are approximately $0.716826684$ and $0.754832189\pF$.
The residual gain and GBW offsets follow from the selected grid-constrained
dimensions. Compliance with the specifications requires simulation
with the corresponding PDK devices and realizable capacitor values.

\newpage
\textbf{Dummy devices and physical implementation.}
Each active MOS has one separately named dummy instance placed
beside it.
Every dummy's D/G/S/B terminals connect to $V_{SS}$, including
the PFET dummies. Dummy $n_f$ and $m$ remain configurable.
These devices are not electrically paralleled with the active MOS,
so their widths are excluded from $W_{\mathrm{eff}}$ in the tables.
A PFET dummy requires its own $V_{SS}$ well/tap separated from
the active PFET's $V_{DD}$ well. Each dummy retains its own integer
$n_f$ and $m$ and legal per-finger width; the initial dummy geometry
follows its associated active MOS. The physical placement must satisfy
the PDK well-spacing and isolation rules.

\textbf{Matching and common-centroid placement.}
The equal-size PMOS mirror M2/M3 uses identical units in a
$2\times2$ AB/BA common-centroid arrangement, preserving equal
centroids and consistent orientation. The unequal-width PMOS buffer
uses equal-centroid interleaving of its driver and load units;
this is not identical-unit matching and does not remove the devices'
intentional width difference or the load's body effect.
NMOS/PMOS pairs use separate well/body regions and symmetric blocks
associated with each device; no cross-type common-centroid equivalence
is assumed. Dummies remain adjacent to their associated active units
and are not consolidated into a separate array.

\textbf{Multiplicity and extracted-device correspondence.}
The PDK \texttt{mSwitch} convention must agree with the schematic's
width, finger-count, and multiplicity interpretation.
The initial schematic uses \texttt{mSwitch=1}. Its correspondence
to grouped extraction for $n_f>1$ requires CDF evaluation and
comparison of schematic and extracted widths and multiplicities.
Decomposition into units for common-centroid placement must preserve
these same electrical quantities.

\textbf{Consistency of simulated and physical dimensions.}
DC and AC verification use the exact grid-constrained geometry,
selected PDK passive values, and intended $n_f$, $m$, and
\texttt{mSwitch} settings. The schematic and layout retain these
values after the specifications are met. Any subsequent change in
dimensions or effective multiplicity requires renewed circuit simulation.
\end{enumerate}

\newpage
\textbf{Model limitations and circuit verification.}

The first-order model determines the initial dimensions, bias voltages,
and compensation values.
The PMOS-only load requires the body-effect correction above;
the longer-channel CMOS designs depend on the accuracy of the
assumed $\lambda(L)$ behavior. Gate capacitance, finite mirror
reference impedance, junction capacitance, and physical passive
models may shift both gain and bandwidth.

\vspace{8pt}
The DC transfer characteristic is obtained by sweeping $V_{in}$
from $-0.9$ to $+0.9\V$. The two intercepts and the input bias
producing $V_{out}=0$ characterize the operating point. AC analysis
about that bias gives the low-frequency gain, relative $-3$ dB
bandwidth, gain-bandwidth product, and any positive-frequency
unity-gain crossing. Dynamic phase lag excludes the DC inversion
and is evaluated at UGF for the high-gain designs and at the bandwidth
for the unity-gain designs. Each sizing change and parasitic extraction
requires a new bias calculation.

\vspace{8pt}
Pre-layout simulation uses legal, grid-constrained MOS dimensions
and PDK passive values. The schematic and layout share the same
device parameters and connectivity, including $R_S$ and $C_m$.
Any dimensional change requires renewed DC and AC analysis.
DRC and LVS establish geometrical legality and schematic/layout
correspondence; simulation with extracted parasitic capacitances
then determines the resulting circuit performance. The analytical
estimates alone do not establish compliance with these requirements.

\vspace{12pt}
\textbf{Sources.}
\begin{enumerate}[label={[\arabic*]},leftmargin=0.3in,
topsep=4pt,itemsep=5pt,parsep=0pt]
\item S. Palermo, \emph{ECEN 474/704 Lab 7: Frequency Response
of Inverting Amplifiers}, course laboratory handout,
Figs.~7-2, 7-5, 7-6, 7-8, and 7-9; Prelab Questions 1--3.
\item Anonymous contributor, \emph{ECEN 474/704 Lab 1: Introduction to Cadence
and MOS Device Characterization}, laboratory report, Table~I:
single-finger $5/0.27\um$ NFET/PFET parameter measurements.
\end{enumerate}

\end{document}
